\scanpage{13}

Je profite de l'occasion pour soulever une question sur la ``cohomologie relative''. Soit $\unsure{q|g}:X\to Y$ un morphisme de topos. Si $F$ est un faisceau abélien (ou un complexe d'iceux) sur $Y$, peut-on définir \unsure{fon\uline{ctoriellem}ent|fo\uline{nctoriellem}ent} en $F$ la cohomologie relative $\mathbb R\Gamma(Y\mathbin{\mathrm{mod}}X,F)$ (de la catégorie dérivée de $\unsure{\mathbb{A}b|\mathit{Ab}}(Y)$ dans celle de $\mathbb{A}b$) ? L'interprétation ``géométrique'' en termes d'opérations sur des $n$-champs de Picard ($n$ ``grand'') suggère que ça doit exister
Mais je ne vois de construction \unsure{évident|évidente} ``\unsure{\textsuperscript{.}|\phantom{\textsuperscript{.}}}à la main'' que dans les deux cas extrêmes :

a) $q$ est ``$(-1)$-acyclique'' i.e. pour tt $F$ sur $Y$, $F\mathord{-}q_{*}q^{*}F$ est injectif (\uline{NB} c'est le cas de $Y_{/\unsure{p|P}}\to Y$ \unsure{si|où} $p\to e_{Y}$ est un \unsure{épimorphisme} - c'est donc le cas de $B_{\unsure{c|e}}\to B_{\unsure{c|G}}$ plus haut)

On prend
% isolated point at source (74,1603), left of display
\margin{\unsure{.|\phantom{.}}}
\[
\mathbb R\Gamma_{\deleted{\unsure{Y|X}}}\bigl(\operatorname{Coker}\bigl(F\to q_{*}\unsure{[|(}\underbrace{e(q^{*}(F)}_{\text{résol. inj.}}))\bigr)[-1]\bigr)
\]
% isolated point at source (202,1634), above-right of archive number
\margin{\unsure{.|\phantom{.}}}
% archive sheet number: 12
